В современных платформах А/Б-эксперименты вынуждены давать ответ за
ограниченное время --- от нескольких дней до нескольких недель, --- тогда
как ключевые бизнес-метрики (удержание, пожизненная ценность пользователя)
формируются на горизонтах, существенно превышающих длительность одного
теста~\cite{kohavi2020trustworthy,kohavi2016pitfalls,xu2015impact}. Прямые
оценки долгосрочного эффекта через продолжительные эксперименты сопровождаются
ростом стоимости, операционных рисков и эффектов выживаемости и потому
плохо встраиваются в быстрые циклы разработки~\cite{kohavi2015focus,
kohavi2016pitfalls,harper2024long}. Альтернативные подходы на основе
суррогатных моделей опираются на модели <<чёрного ящика>> и требуют
постоянной перекалибровки при изменениях продукта~\cite{tripuraneni2024,
duan2021online,dimakopoulou2023evaluating}. В результате между быстрыми
экспериментальными циклами и стратегической потребностью смотреть на
долгосрочные эффекты сохраняется устойчивый разрыв.

Чтобы конкретизировать характер этого разрыва, рассмотрим реальный пример
из практики. В рассматриваемом эксперименте пользователям тестовой группы
предлагалась возможность <<купить сейчас, заплатить позже>> (Buy Now,
Pay Later, BNPL) в маркетплейсе. На рис.~\ref{fig:retention_diff} показано,
как менялась разница в показателях удержания пользователей контрольной (A)
и тестовой (B) групп.

\begin{comment}
\begin{figure}[h]
\centering
\includegraphics[width=1.0\linewidth]{data/01_retention_diff_ci.png}
\caption{Разница в показателях удержания (тестовая группа~B минус
контрольная группа~A) в эксперименте BNPL. По оси~$y$ --- абсолютная разница
ежедневных ставок удержания, по оси~$x$ --- число дней с момента старта
эксперимента. Устойчиво положительная разница, начиная с 41-го дня (момент
окончания эксперимента), указывает на сохраняющееся улучшение удержания
в тестовой группе.}
\label{fig:retention_diff}
\end{figure}
\end{comment}

На старте эксперимента уровни удержания в обеих группах были близки. По ходу
эксперимента BNPL воспринималась пользователями как ценная функциональность,
что приводило к росту удержания в тестовой группе. К моменту завершения
эксперимента сформировавшийся разрыв сохранялся в течение длительного периода.
Этот пример показывает, что долгосрочные поведенческие паттерны могут
меняться \textit{в течение} эксперимента и что соответствующие эффекты
допустимо оценивать ещё до его формального окончания: тестовая группа
ожидаемо приносит большую ценность независимо от итогового решения по самой
функциональности.

Опираясь на этот пример, мы интерпретируем долгосрочный эффект не как ответ
на гипотетический вопрос <<что произошло бы, если бы эксперимент длился
бесконечно?>>, а более операциональным образом: <<как именно эксперимент
уже изменил долгосрочное поведение пользователей?>>. Такая постановка
позволяет работать с фактически наблюдаемыми поведенческими сдвигами и
оценивать их до того, как реализуется вся будущая выручка, --- именно эта
интерпретация лежит в основе предлагаемого метода Orbitals.

Неформально задача исследования состоит в следующем. Дана платформа,
на которой проводится А/Б-эксперимент с разделением пользователей на
контрольную (A) и тестовую (B) группы; известны их характеристики и
история активности до и в течение эксперимента. Требуется построить
\textit{целеориентированную} поведенческую сегментацию пользователей
(орбитали) и метод оценки долгосрочного эффекта эксперимента, который
удовлетворяет следующим требованиям:
\begin{enumerate}[label=(\roman*),leftmargin=*,itemsep=2pt]
  \item использует только данные, наблюдаемые в течение самого теста;
  \item опирается на изменение распределения пользователей по орбиталям
        и переходы между ними;
  \item сохраняет интерпретируемость относительно исходных поведенческих
        сегментов;
  \item встраивается в существующие пайплайны А/Б-тестирования без
        существенной перестройки инфраструктуры.
\end{enumerate}

\begin{comment}
\textbf{Объектом исследования} являются поведенческие паттерны пользователей
онлайн-платформы и их изменение под воздействием экспериментальных
интервенций. \textbf{Предметом исследования} --- методы целеориентированной
сегментации и моделирования траекторий пользователей для оценки долгосрочного
эффекта А/Б-экспериментов.

\textbf{Цель} магистерской работы --- разработка и исследование метода
Orbitals для оценки долгосрочного эффекта онлайн-экспериментов на основе
динамики поведенческой сегментации пользователей. Для достижения цели
поставлены следующие \textbf{задачи}:
\begin{itemize}
  \item проанализировать и систематизировать существующие подходы к оценке
        долгосрочного эффекта онлайн-экспериментов и моделированию
        жизненного цикла пользователей;
  \item формализовать постановку задачи в терминах ограниченной LTV и
        целеориентированной поведенческой сегментации;
  \item разработать метод построения орбиталей и оценки переходов между
        ними в А/Б-экспериментах;
  \item предложить метрики, позволяющие по изменению распределения
        пользователей по орбиталям оценивать долгосрочный uplift;
  \item экспериментально исследовать метод Orbitals на данных
        крупномасштабной онлайн-платформы и сопоставить его с базовыми
        и суррогатными подходами.
\end{itemize}

\textbf{Научная новизна} работы состоит в формализации концепции орбиталей
как целеориентированных поведенческих сегментов и в построении на их основе
интерпретируемого метода оценки долгосрочного эффекта по короткому
А/Б-эксперименту, замыкающего разрыв между суррогатным моделированием и
динамической сегментацией клиентской базы.

\textbf{Практическая значимость} заключается в том, что предложенный метод
может быть встроен в действующие платформы А/Б-тестирования крупных
онлайн-сервисов и использоваться для более ранней и устойчивой оценки
долгосрочного эффекта продуктовых изменений.

Формальная математическая модель, лежащая в основе предлагаемого метода ---
определение ограниченной LTV, орбиталей и структуры переходов между ними,
а также вытекающие из неё метрики эффекта, --- подробно описывается в
главах, посвящённых обзору литературы и описанию метода Orbitals.
\end{comment}